% Generated by scripts/materialize_unified_paper.py; do not edit.
\section{Gaussian source ledger and notation}\label{low:sec:setup}

Fix \(C>1\).  Let \(p_C\in(0,1/2)\) be the solution of
\[
 \log p_C=C\log(1-p_C),
\]
and, in the range of \(C\) used below, set
\[
 p=p_C+\frac1C<\frac12,\qquad
 \Phi(-c)=p,\qquad a=\varphi(c),\qquad
 D=\frac{4aC}{1-p}.
\tag{L2.1}\label{low:main:eq:2-1}
\]
For the fixed-\(C\) limit \(\ell\to\infty\), take
\[
 d=\left\lceil D^2\ell^2\right\rceil,\qquad
 D_\ell=\frac{\sqrt d}{\ell}=D+O_C(\ell^{-2}).
\tag{L2.2}\label{low:main:eq:2-2}
\]
The harmless replacement of \(D\) by \(D_\ell\ge D\) in finite-\(\ell\)
feasibility inequalities will always be left implicit.

For a standard normal variable conditioned on \(Z\le-b\), define
\[
 m(b)=\frac{\varphi(b)}{\Phi(-b)},\qquad
 v(b)=1+b\,m(b)-m(b)^2,\qquad Y_b=Z+m(b),
\tag{L2.3}\label{low:main:eq:2-3}
\]
so that \(Y_b\) is centered and has variance \(v(b)\).  Its
positive-direction cumulant generating function is
\[
 K_b(u)=\log\mathbb E e^{uY_b}
 =\frac{u^2}{2}+u\,m(b)+\log\Phi(-b-u)-\log\Phi(-b),
 \qquad u\ge0,
\tag{L2.4}\label{low:main:eq:2-4}
\]
and we write
\[
 \psi_b(u)=
 \begin{cases}K_b(u)/u^2,&u>0,\\[1mm]v(b)/2,&u=0.
 \end{cases}
\tag{L2.5}\label{low:main:eq:2-5}
\]
We shall repeatedly use
\[
 K_b''(u)=v(b+u),\qquad
 K_b(u)=\int_0^u(u-s)v(b+s)\,ds.
\tag{L2.6}\label{low:main:eq:2-6}
\]

\begin{lemma}[One-sided Gaussian calculus]\label{low:lem:mills}
For every \(b\in\mathbb R\), \(m'(b)=m(b)(m(b)-b)>0\),
\(m''(b)\ge0\), and \(v(b)=1-m'(b)\) is nonincreasing.  Consequently,
\(0<v(b)\le1\), \(K_b''(u)\in[0,1]\) for \(u\ge0\), and
\(\psi_b(u)\) is nonincreasing in both \(b\) and \(u\).
\end{lemma}

\begin{proof}
Differentiating \(m=\varphi/\Phi(-\cdot)\) gives
\(m'=m(m-b)\).  The classical inverse-Mills inequality \(m(b)>b\) for
\(b>0\) (and positivity of \(m\) for \(b\le0\)) gives \(m'>0\).
Moreover \(v=1-m'\).  Monotonicity of the variance of an upper-truncated
normal in its upper cutoff, applied at the cutoff \(-b\), says that \(v\) is
nonincreasing in \(b\); equivalently \(m''=-v'\ge0\).  The variance bounds
give \(0<v\le1\).  Formula \eqref{low:main:eq:2-6} then implies the assertions concerning
\(K_b\) and \(\psi_b\); for example
\(\psi_b(u)=\int_0^1(1-t)v(b+ut)\,dt\).
\end{proof}

Put
\[
\begin{aligned}
 m_0&=m(c)=\frac ap,&m'_0&=m'(c),\\
 \alpha_R&=10\sqrt{\log(10/p)},&
 \omega_0&=\frac{\alpha_R^2}{D}
 =\frac{100\log(10/p)}D .
\end{aligned}
\tag{L2.7}\label{low:main:eq:2-7}
\]
and
\[
 \rho=\frac{m_0m'_0}{D},\qquad
 \theta=m_0^2,\qquad A_0=m_0^2m'_0.
\tag{L2.8}\label{low:main:eq:2-8}
\]
For fixed \(w>\omega_0\), define the lower and upper cutoff extrema
\[
\begin{aligned}
 m_w&=m(c-w),&M_w&=m(c+w),\\
 j_-(w)&=m'(c-w),&j_+(w)&=m'(c+w).
\end{aligned}
\tag{L2.9}\label{low:main:eq:2-9}
\]
The distinction between \(m_w\) and \(M_w\) is essential: the former gives
the history-uniform exact-CGF deficit, whereas the latter bounds the new
linear tilts.

\subsection{The frozen red and blue source ledgers}

Fix one absolute source constant \(K\) that dominates the red
\(O(\sqrt{\log(1/p)}/D)\) errors in the pinned HMS induction, and put
\[
 \vartheta_C=1-K\frac{\sqrt{\log(1/p)}}D>0,
\tag{L2.10}\label{low:main:eq:2-10}
\]
\[
 B_R(C)=-\log\frac p{p_C}
 +\vartheta_C\frac{a^3}{3p^3D},\qquad
 B_B(C)=\log\frac{1-p_C}{1-p}
 -\frac{4a^3C}{3(1-p)^3D}.
\tag{L2.11}\label{low:main:eq:2-11}
\]
The literal red target envelope is
\[
 \mathcal H_C(r)=p^{\binom r2}
 \exp\!\left[-\vartheta_C\frac{a^3}{p^3\sqrt d}\binom r3
 +\frac apd^{-1/5}\binom r2+r\right].
\tag{L2.12}\label{low:main:eq:2-12}
\]
Thus \(B_R/2\) and \(CB_B/2\), respectively, are the exact limiting red
and blue source bonuses relative to \(-\tfrac12\log p_C\).  In particular,
these symbols do not denote the conclusion of the present paper.

\subsection{The pre-existing exact-CGF deficit}

For \(1\le k\le\ell\), set
\[
 P^0_{k,d}=\frac{\sqrt d}{\sqrt d-4m_0k},
\tag{L2.13}\label{low:main:eq:2-13}
\]
and, whenever the denominator is positive, define
\[
 \mathsf d_{k,w,d}=
 \frac{m_0^2m_w^2}{d}k(k-1)^2
 \left[1-P^0_{k,d}\psi_{c-w}\!\left(
 \frac{P^0_{k,d}m_0m_w(k-1)}{\sqrt d}\right)\right],
\tag{L2.14}\label{low:main:eq:2-14}
\]
\[
 \mathcal D_{C,w,d}(r)=\sum_{k=1}^{r-1}\mathsf d_{k,w,d}.
\tag{L2.15}\label{low:main:eq:2-15}
\]
For \(0\le x\le1\), let
\[
 P_*(x)=\frac D{D-4m_0x},\qquad
 u_*(x)=\frac{P_*(x)m_0m(c-\omega_0)x}{D}.
\tag{L2.16}\label{low:main:eq:2-16}
\]
For all sufficiently large \(C\), \(D>8m_0\) and
\(P_*(1)v(c-\omega_0)<2\); hence \eqref{low:main:eq:2-13}--\eqref{low:main:eq:2-17} have positive
denominators and nonnegative integrands throughout the range used below.
The exact source-relative CGF term used in the theorem is
\[
 G_*(C)=\frac{m_0^2m(c-\omega_0)^2}{D^2}
 \int_0^1x^3\left[1-P_*(x)
 \psi_{c-\omega_0}\bigl(u_*(x)\bigr)\right]dx.
\tag{L2.17}\label{low:main:eq:2-17}
\]
Equivalently, if \(G(C,w)\) is obtained from \eqref{low:main:eq:2-17} by replacing
\(\omega_0\) with \(w\), then
\[
 \lim_{\ell\to\infty}\ell^{-2}\mathcal D_{C,w,d}(\ell)=G(C,w),
 \qquad G_*(C)=\lim_{w\downarrow\omega_0}G(C,w).
\tag{L2.18}\label{low:main:eq:2-18}
\]

\subsection{The controlled-residual constants}

Set
\[
 a_- =\min\{A_0,\theta\},\qquad
 a_+=\max\left\{\theta,
 \frac{A_0+\rho\theta}{1-\rho}\right\},
\tag{L2.19}\label{low:main:eq:2-19}
\]
\[
 \mu_w=j_-(w)-j_+(w)^2
 \frac{m_0+\theta/D}{D(1-\rho)},
\tag{L2.20}\label{low:eq:mu-w}
\]
and
\[
 \widehat L_w(C)=
 \frac{4m_0M_w^2}{D}
 +\frac{2M_w^2a_+}{D^2}
 +\frac{32m_0^2M_w^2a_+}{a_-D^2}.
\tag{L2.21}\label{low:main:eq:2-21}
\]
Finally put
\[
 \widehat B_w(C)=m_0^2-\frac{16}{3}-\widehat L_w(C),
 \qquad
 E_w(C)=\widehat B_w(C)\theta-\frac{\theta^2}{2\mu_w},
\tag{L2.22}\label{low:main:eq:2-22}
\]
\[
 \widehat H(C,w)=
 \frac{\widehat B_w(C)A_0+E_w(C)}{12D^2},
 \qquad
 \widehat H_*(C)=\lim_{w\downarrow\omega_0}\widehat H(C,w).
\tag{L2.23}\label{low:main:eq:2-23}
\]

\subsection{Probability space, exposure histories, and source events}

For each \(0\le r\le\ell\), let
\(\boldsymbol x_1,\ldots,\boldsymbol x_r\) be independent
\(N(0,d^{-1}I_d)\) vectors.  An edge is blue when
\(\langle\boldsymbol x_i,\boldsymbol x_j\rangle\ge-c/\sqrt d\) and red
otherwise.  We use the Bartlett representation
\(\boldsymbol y_1,\ldots,\boldsymbol y_r\): the entries
\(y_j(i)\), \(i<j\), are independent \(N(0,d^{-1})\), the diagonal
entries are independent with
\(y_i(i)\sim\sqrt{\chi^2_{d-i+1}/d}\), and \(y_j(i)=0\) for \(i>j\).
This representation has the same joint Gram matrix as the original Gaussian
vectors.

Let \(M_i=(y_j(i))_{j=i+1}^r\) be the strict part of column \(i\), and put
\[
 \mathcal D_r=\sigma\bigl(y_1(1),\ldots,y_r(r)\bigr).
\]
Define the exposure filtration
\[
 \mathcal F_s=\mathcal D_r\vee\sigma(M_1,\ldots,M_s),
 \qquad 0\le s\le r.
\tag{L2.24}\label{low:main:eq:2-24}
\]
All conditional assertions below are statements about regular conditional
probabilities and hold almost surely.  Write \(\pi_s\) for projection onto
the first \(s\) Bartlett coordinates and put
\[
 h_{ij}=\langle\pi_{i-1}(\boldsymbol y_i),
                   \pi_{i-1}(\boldsymbol y_j)\rangle,
 \qquad 1\le i<j\le r.
\tag{L2.25}\label{low:main:eq:2-25}
\]
Thus \(h_{ij}\) is \(\mathcal F_{i-1}\)-measurable.  Let \(E_{ij}\) be the
blue-edge event and set
\[
 I_s=\bigcap_{s<i<j\le r}\overline E_{ij},\qquad
 C_s=\bigcap_{s<i<j\le r}E_{ij}.
\tag{L2.26}\label{low:main:eq:2-26}
\]

For the source perfectness events put
\(\alpha_B=10\sqrt{C\log(10/p)}\) and
\(\delta=\alpha_Bd^{-1/4}\).  The event \(B^R_r\) says that, for every
\(i\le r\),
\[
 \|\boldsymbol y_i\|\in(1-\delta,1+\delta),\qquad
 \|\pi_{i-1}(\boldsymbol y_i)\|
 \le \frac{\alpha_R\sqrt\ell}{\sqrt d};
\tag{L2.27}\label{low:main:eq:2-27}
\]
\(B^B_r\) is defined by replacing \(\alpha_R\) with \(\alpha_B\).  We use
the probability notation
\[
 P^*_{R,r}=\mathbb P(I_0\cap B^R_r),\quad
 P_{R,r}=\mathbb P(I_0),\quad
 P^*_{B,r}=\mathbb P(C_0\cap B^B_r),\quad
 P_{B,r}=\mathbb P(C_0).
\tag{L2.28}\label{low:main:eq:2-28}
\]
Empty intersections have probability one.

For \(0\le s\le r\), let
\[
 \kappa^R_{r,s}=\mathbb P(B^R_r\mid\mathcal F_s),\qquad
 \mathfrak H^R_{r,s}=\{\kappa^R_{r,s}>0\}.
\tag{L2.29}\label{low:main:eq:2-29}
\]
We call an \(\mathcal F_s\)-history \emph{red-admissible} when it belongs to
\(\mathfrak H^R_{r,s}\).  This definition is independent of a choice of
coordinates for the conditional kernel.  On the complementary histories,
\(\mathbb P(I_s\cap B^R_r\mid\mathcal F_s)=0\) almost surely.  Thus every
conditional estimate below only needs to be checked on
\(\mathfrak H^R_{r,s}\).

At exposure \(i\), let
\(A_i^R=\bigcap_{j>i}\overline E_{ij}\).  Conditional on
\(A_i^R\) and \(\mathcal F_{i-1}\), set \(X_j=y_j(i)\), \(j>i\).  Its exact
upper cutoff is
\[
 X_j\le-\frac{b_{ij}}{\sqrt d},\qquad
 b_{ij}=\frac{c+\sqrt d\,h_{ij}}{y_i(i)}.
\tag{L2.30}\label{low:main:eq:2-30}
\]

The red-perfect extraction map \(\Gamma_R\) scans
\(\boldsymbol x_1,\ldots,\boldsymbol x_\ell\) in order and retains an index
when adjoining its vector preserves the two inequalities in \eqref{low:main:eq:2-27}.  For a
specified \(I\subseteq[\ell]\), define
\[
 p_I=\mathbb P\bigl(I_0\cap\{\Gamma_R=I\}\bigr).
\tag{L2.31}\label{low:main:eq:2-31}
\]
This fixes both the event and the retained-order relabeling used later.

For a deterministic nonnegative triangular array
\(T=(T_{jq})_{1\le j<q\le r}\), define the weighted projection potential
\[
 \mathcal P_s^T=
 \sum_{s<j<q\le r}T_{jq}
 \langle\pi_s(\boldsymbol y_j),\pi_s(\boldsymbol y_q)\rangle.
\tag{L2.32}\label{low:main:eq:2-32}
\]
For \(i\ge1\), its normalized one-step kernel is
\[
\begin{aligned}
 \mathcal K_i^T={}&
 \exp\!\left(\sum_{j>i}T_{ij}h_{ij}\right)
 \mathbb P(A_i^R\mid\mathcal F_{i-1})\\
 &\times\mathbb E\!\left[
 \exp\!\left\{-\sum_{i<j<q\le r}T_{jq}X_jX_q\right\}
 \middle|A_i^R,\mathcal F_{i-1}\right].
\end{aligned}
\tag{L2.33}\label{low:main:eq:2-33}
\]

\begin{assumption}[Version-pinned source and weighted-extension interface]
\label{low:ass:hms}
Items (S1), (S2), (S4), and (S5) are imported from the content-addressed
HMS v2 and Lin--Niu v2 sources identified in Appendix~\ref{low:app:bridge}.
Item (S3) is the exact weighted-extension hypothesis required by this paper;
it is not represented as a theorem stated verbatim in either source.
\begin{enumerate}[label=(S\arabic*)]
\item For every \(r\le\ell\) and \(1\le i<r\), on every red-admissible
\(\mathcal F_{i-1}\)-history, the next-column coordinates in \eqref{low:main:eq:2-30},
conditionally on \(A_i^R\), are independent \(N(0,d^{-1})\) variables
truncated by the displayed upper cutoffs, and
\[
 |b_{ij}-c-\sqrt d\,h_{ij}|\le d^{-1/5},\qquad
 |\sqrt d\,h_{ij}|\le\omega_0.
\tag{L2.34}\label{low:main:eq:2-34}
\]
Thus, for every fixed \(w>\omega_0\), all cutoffs lie in
\([c-w,c+w]\) once \(\ell\ge\ell_0(C,w)\).
\item Under the conditioning and history quantifiers in (S1), with
\(\zeta_j=\sqrt d(X_j-\mathbb E[X_j\mid
A_i^R,\mathcal F_{i-1}])\), the variables \((\zeta_j)_{j>i}\) are
conditionally independent, centered, and \(1\)-subgaussian, and
\[
 \log\mathbb E[e^{x\zeta_j^2}\mid A_i^R,\mathcal F_{i-1}]
 \le\frac{16}{3}x,\qquad 0\le x\le\frac18.
\tag{L2.35}\label{low:main:eq:2-35}
\]
For \(k=r-i\), \(Q_0=d/(4t_0k)\), and centered \(\zeta_j\), the frozen
uniform-quadratic factor obeys
\[
 \frac1{Q_0}\log\mathbb E\exp\!\left\{
 -\frac{Q_0t_0}{d}\sum_{i<j<q\le r}\zeta_j\zeta_q\right\}
 \le\frac{4t_0k}{d}.
\tag{L2.36}\label{low:main:eq:2-36}
\]
The positive-direction CGF is exactly \eqref{low:main:eq:2-4}.  In the source pre-Cauchy
allocation, its centered-linear proxy is precisely
\(\sum_{j>i}(u_j^0)^2\), with \(u_j^0\) as in \eqref{low:main:eq:5-9}; replacing that proxy
by the exact CGF produces the same-history deficit \eqref{low:main:eq:5-16}.
\item For every target \(r\le\ell\), every deterministic nonnegative
triangular \(T\) fixed before exposure, and every deterministic real sequence
\(\lambda_1,\ldots,\lambda_{r-1}\), the following weighted reverse-induction
rule holds.  If, on every history in \(\mathfrak H^R_{r,i-1}\), with
\(k=r-i\),
\[
 \mathcal K_i^T\le
 \frac{\mathcal H_C(k+1)}{\mathcal H_C(k)}e^{-\lambda_i},
 \qquad 1\le i<r,
\tag{L2.37}\label{low:main:eq:2-37}
\]
then, for every \(0\le s\le r\),
\[
 \mathbb P(I_s\cap B^R_r\mid\mathcal F_s)
 \le\mathcal H_C(r-s)
 \exp\!\left[-\mathcal P_s^T-
 \sum_{i=s+1}^{r-1}\lambda_i\right].
\tag{L2.38}\label{low:main:eq:2-38}
\]
On \((\mathfrak H^R_{r,s})^c\), the left side of \eqref{low:main:eq:2-38} is zero, so the same
inequality holds almost surely on the entire conditional state space.
The source law of total probability is
\[
 \mathbb P(I_{i-1}\cap B^R_r\mid\mathcal F_{i-1})
 =\mathbb P(A_i^R\mid\mathcal F_{i-1})
 \mathbb E[\mathbb P(I_i\cap B^R_r\mid\mathcal F_i)
 \mid A_i^R,\mathcal F_{i-1}],
\tag{L2.39}\label{low:main:eq:2-39}
\]
and the exact projection identities are
\[
\begin{aligned}
 \mathcal P_i^T
 &=\sum_{i<j<q\le r}T_{jq}
 \langle\pi_{i-1}\boldsymbol y_j,\pi_{i-1}\boldsymbol y_q\rangle
   +\sum_{i<j<q\le r}T_{jq}X_jX_q,\\
 \mathcal P_{i-1}^T
 &=\sum_{i<j<q\le r}T_{jq}
 \langle\pi_{i-1}\boldsymbol y_j,\pi_{i-1}\boldsymbol y_q\rangle
   +\sum_{j>i}T_{ij}h_{ij}.
\end{aligned}
\tag{L2.40}\label{low:main:eq:2-40}
\]
Equations \eqref{low:main:eq:2-33}, \eqref{low:main:eq:2-39}, and \eqref{low:main:eq:2-40} prove \eqref{low:main:eq:2-38} by reverse induction
from \(s=r\).  For the uniform array \(T^0_{jq}=t_0\), the version-pinned
source scalar allocation, together with the exact-CGF replacement in (S2),
is the explicit baseline statement on every history in
\(\mathfrak H^R_{r,i-1}\):
\[
 \mathcal K_i^{T^0}\le
 \frac{\mathcal H_C(k+1)}{\mathcal H_C(k)}
 e^{-\mathsf d_{k,w,d}}
 \qquad(k=r-i).
\tag{L2.41}\label{low:main:eq:2-41}
\]
Equation \eqref{low:main:eq:2-41}, rather than a final Ramsey bound, is the scalar source input.
No nonuniform one-step estimate is imported.  Sections
\ref{low:sec:residual}--\ref{low:sec:induction} compare the nonuniform kernel with
\eqref{low:main:eq:2-41} and prove \eqref{low:main:eq:2-37} for the array used here.
\item If red-perfect extraction deletes \(u\) vertices from a target of size
\(\ell\), then, for each specified retained set \(I\) of size \(\ell-u\),
\[
 p_I\le P^*_{R,\ell-u}\left(\frac p{10}\right)^{10\ell u},
\tag{L2.42}\label{low:main:eq:2-42}
\]
for every \(0\le u\le\ell\), with
\(P^*_{R,0}=\mathcal H_C(0)=1\).  The probability \(P^*_{R,\ell-u}\)
uses the retained-order relabeling defined after \eqref{low:main:eq:2-31}.
\item For each fixed \(C\), there is a real sequence
\(\epsilon_{C,\ell}\to0\) such that, with
\(q_\ell=\lfloor C\ell\rfloor\), the unchanged blue induction and its
perfect-sequence extraction give
\[
 P_{B,q_\ell}\le(1-p)^{\binom{q_\ell}{2}}
 \exp\left[
 \frac{4a^3}{(1-p)^3\sqrt d}\binom{q_\ell}{3}
 +\epsilon_{C,\ell}\ell^2\right].
\tag{L2.43}\label{low:main:eq:2-43}
\]
\end{enumerate}
No Ramsey-rate conclusion, no red--blue comparison, and no new residual term
is included in (S1), (S2), (S4), or (S5).  The only additional source-side
hypothesis is the precisely delimited weighted rule (S3).
\end{assumption}

The constant \(K\) in \eqref{low:main:eq:2-10} is not effective in the pinned source.
Consequently all sufficiently-large-\(C\) statements below are existential.
